% ============================================================
% Part III — Desktop Section — APPENDIX: Friction Catalogue — DRAFT v1
% Ordered 2026-07-24 (user directive, M3.1 review): every reported friction,
%   most to least reported. Referenced by the M3.1 ¶2 footnote.
% Data: A1 friction family, all-Fable corpus (refreshed 2026-07-23);
%   counts = coded applications (t3-stats); percentages = unique STUDENTS
%   mentioning the friction, of app users n=123 / YouTube group n=108
%   (strata from Q8; computed 2026-07-24 from 03-coded-corpus × respondents).
% Exemplars: curated bank (04-exemplar-bank) or verbatim evidence spans (03).
% Residual grouping (A1.other): DESCRIPTIVE grouping for presentation only —
%   not codes; the 32 underlying applications are listed in 03 under A1.other.
% Verify all totals against the corpus at assembly (as coding appendix).
% Appendix number + \label at assembly.
% ============================================================

\section*{Appendix: Catalogue of Reported Frictions}

Every friction students reported is collected here, ordered from most to
least reported. A complaint is counted once per coded application under the
friction family of the codebook; the percentage columns count students, not
complaints---the share of the 123 application users, and of the 108 students
on the YouTube route, who raised the friction at least once. Two of the
frictions we predicted during development appear at the bottom of the table
precisely because the survey did not bear them out.

\begin{table}[h]
\centering
\small
\begin{tabular}{lccl}
\hline
Friction & Complaints & App users / YouTube & Example from the survey \\
\hline
Crashes and freezes & 29 & 15\% / 8\% &
  ``I had issues with crashes'' (W25-R064/Q6) \\
Install and compatibility & 26 & 3\% / 17\% &
  ``[m]ore inclusive for macbook users :('' (W25-R042/Q9) \\
Buffering and latency & 25 & 11\% / 8\% &
  ``maybe improving the lags, I would glitch sometimes'' (W25-R095/Q9) \\
Wayfinding and navigation & 21 & 11\% / 6\% &
  ``eliminate having to go up stairs and stuff?'' (W25-R084/Q9) \\
Hardware access & 13 & 3\% / 7\% &
  ``the file being too large, my computer could not handle it'' (W25-R060/Q9) \\
Video fidelity & 10 & 5\% / 3\% &
  ``I feel like it's very blurry in my experience'' (W25-R062/Q9) \\
Viewing proximity & 2 & --- &
  ``surgical screen far in the back which made it hard to see'' (S25-R031/Q9) \\
Avatar occlusion & 0 & --- &
  (predicted; unattested---most students entered alone) \\
\hline
\end{tabular}
\caption{Reported frictions by frequency. Counts are coded applications;
percentages are students mentioning the friction at least once.}
\end{table}

A further 32 complaints name frictions outside these categories and are
carried in the codebook as a residual group. Grouped descriptively, they
comprise: difficulty understanding the footage or the task it served (5,
e.g.\ ``hard to tell what is happening in them,'' W25-R067/Q6); interface
usability (5, e.g.\ ``The VR platform is not very user friendly,''
W26-R039/Q9); distraction and overstimulation (4, e.g.\ ``more distracting
than anything,'' W25-R092/Q9); motion discomfort and disorientation (4,
e.g.\ ``I get motion sickness when watching the VR videos,'' W25-R054/Q9);
residual crash, freeze, and loading complaints phrased outside the crash
category (4); missing or inaudible audio on the surgical footage (3, e.g.\
``Most surgery videos are silent,'' W25-R039/Q9); passivity and filler time
(2, e.g.\ ``a lot of filler space/period of inactivity,'' W25-R025/Q9); and
five single mentions: wayfinding signage, account security, solitude in the
platform (``I was usually alone in the platform,'' W25-R038/Q10), setup
abandonment, and general doubt of benefit. Not counted as frictions are the
19 students on the YouTube route (18\% of that group) who simply reported
never using the platform at all.
